% manuscript/drafts/introduction.tex
%
% Session L4, Step 13 draft. Session L11 refines and integrates this into the
% Springer Nature manuscript. Substantively complete, not submission-ready.
%
% Every factual claim traces to docs/literature/CLAIM_MAP.md. Every citation key
% resolves to manuscript/references.bib, whose entries all carry a DOI verified
% against the Crossref work record during Session L4.
%
% Physics-first framing (docs/CONVENTIONS.md): mechanism and regime are the
% through-line; the mathematics serves the physical statement, never replaces it.

\section{Introduction}
\label{sec:intro}

\subsection{One conservation law, two regimes}

A column of fluid on a rotating planet carries a label it cannot change. Its
potential vorticity,
%
\begin{equation}
q = \frac{\zeta + f}{h},
\label{eq:intro-pv}
\end{equation}
%
combining the column's own spin $\zeta$, the spin $f = 2\Omega\sin\varphi$ it has
merely by sitting on a rotating sphere, and its depth $h$, is materially
conserved: $\mathrm{D}q/\mathrm{D}t = 0$. Two older statements, divided into one
another, give this immediately. Kelvin's circulation theorem fixes
$(\zeta+f)\,\delta A$ for a material column of cross-section $\delta A$; mass
conservation fixes $h\,\delta A$; the unknown area cancels.

Almost everything in the large-scale atmosphere that this paper is concerned with
follows from that single statement, and it follows in two regimes that are usually
treated separately.

Push a column poleward and the planetary vorticity it feels increases. If the free
surface is stiff, so the column cannot stretch, its relative vorticity must fall
by exactly as much. A chain of such displacements therefore acquires an
alternating vorticity pattern whose inverted flow field sits a quarter wavelength
out of phase with the displacement itself --- and a quarter-phase forcing
translates a pattern rather than amplifying it. The disturbance propagates, and it
propagates westward, unconditionally. This is the Rossby wave, and on a sphere its
angular phase speed takes the closed form $c_{\mathrm{ang}} = -2\Omega/[n(n+1)]$
for a mode of total degree $n$.

Now curve a zonal jet sharply enough that its own vorticity gradient locally
cancels the planetary one. The background potential vorticity is no longer
monotonic in latitude, a displaced column no longer always finds more $q$ ahead of
it, and the restoring mechanism loses its single sign. What was a wave becomes an
instability, drawing energy from the shear rather than riding on it. The formal
statement of the loss is the Rayleigh--Kuo criterion \citep{kuo1949}: a growing
normal mode requires the background potential-vorticity gradient to change sign
somewhere in the domain.

These are not two subjects. \citet{bretherton1966} showed that wave development in
a flow with two distinct potential-vorticity gradients is the interaction of two
Rossby waves propagating on them, an idea named and developed as
counter-propagating Rossby waves by \citet{hoskins1985} and given quantitative
form for the barotropic shear layer by \citet{heifetz1999}. In that picture the
wave and the instability are the same mechanism seen at one potential-vorticity
interface and at two: a single interface supports an edge wave that propagates; two
interfaces of opposite sign each sit in the other's induced flow, and if they can
hold a fixed relative phase against the shear between them, each amplifies the
other. This paper is organised around that identification, which we adopt as
exposition rather than claim as new.

\subsection{What is actually hard about testing this}

The theory above is old and secure. What is not secure, and what this paper
addresses, is the quality of the numerical evidence usually brought to bear on it.

Three difficulties recur.

First, the analytic target is only asymptotically correct. The Rossby--Haurwitz
speed $-2\Omega/[n(n+1)]$ is the nondivergent result. A real shallow-water layer
has a finite deformation radius $L_{\mathrm d} = \sqrt{gH}/2\Omega$, so a column
can stretch, and part of the potential-vorticity budget that would have gone into
relative vorticity goes into stretching instead. The restoring agent is diluted
and the wave slows. The relevant parameter is Lamb's parameter
$\varepsilon = 4\Omega^2R^2/(gH) = (R/L_{\mathrm d})^2$, and for Earth-like values
$\varepsilon \approx 8.8$: the planet is about three deformation radii across, so
this is not a small correction. The eigenstructure of the divergent problem is
classical --- it is the content of Laplace's tidal equations, whose
eigenfunctions were computed by \citet{longuethiggins1968}, treated in vector
harmonic form by \citet{swarztrauber1985}, and reviewed with the three-family
branch structure made explicit by \citet{kasahara1976} --- but a solver that
compares itself against the nondivergent formula while integrating the divergent
system is comparing against the wrong number.

Second, the instability criterion is one-sided. Rayleigh--Kuo forbids growth when
the potential-vorticity gradient does not change sign, but when it does change
sign the criterion says nothing about whether anything grows, how fast, or at
which zonal wavenumber. \citet{kuo1949} says so himself: dynamic stability
analysis ``can only give some indication of the possibility for the development of
some disturbances.'' Converting the prohibition into a prediction requires solving
an eigenvalue problem, which is standard --- on the plane since
\citet{kuo1949}, on the sphere in the sustained programme of
\citet{skiba2004,skiba2008,skiba2024} and in the rigorous treatments of
\citet{constantin2022}.

Third, the canonical test cases are less innocent than they look.
\citet{williamson1992} proposed the standard shallow-water suite that the field
still uses, with reference solutions from \citet{jakobchien1995}, and its
zonal-wavenumber-4 Rossby--Haurwitz case was chosen because that wave was believed
stable. \citet{thuburn2000} showed it is not: the shallow-water
wavenumber-4 Rossby--Haurwitz wave is dynamically unstable, breaks down if
perturbed, and --- unlike the nondivergent case --- develops a potential-enstrophy
cascade to small scales, with the outcome depending on the dissipation
coefficient. \citet{galewsky2004} made the parallel point about the suite as a
whole, that its later cases carry ``unexpected subtleties which render them less
useful than originally thought'', and proposed a barotropically unstable jet as a
replacement. A study that initialises Rossby--Haurwitz modes and measures their
phase speed must therefore state its measurement window, not assume a steady
solution.

\subsection{Contributions}

This paper does not contribute new theory. It contributes a controlled,
separately verified numerical isolation of the beta effect as a single independent
variable, spanning the range from linear wave propagation to nonlinear shear
instability in one model and one framework, with quantified numerical uncertainty
and an explicit closure against reanalysis. Specifically:

\begin{enumerate}
\item \textbf{One mechanism carried across two regimes with one solver.} The wave
      results and the instability results are produced by the same code, the same
      conventions and the same uncertainty budget, so they are commensurable. The
      literature treats the two regimes in largely separate bodies of work; the
      contribution here is integration, not discovery.

\item \textbf{Rotation rate as a genuine free parameter.} $\Omega$ is swept over a
      sixteen-fold range, $0.25\,\Omega_0$ to $4\,\Omega_0$, turning a fixed
      physical constant into an independent variable. This allows
      $c_{\mathrm{ang}} \propto \Omega$ for the wave, $L_{\mathrm R} \propto
      \Omega^{-1/2}$ for the Rhines scale \citep{rhines1975,vallis1993}, and the
      stabilisation of a fixed jet by increasing $\Omega$ to be tested rather than
      assumed.

\item \textbf{Verification reported separately from validation.} Convergence
      against analytic solutions \citep{williamson1992,lauter2005}, discretisation
      error from a resolution ladder, hyperdiffusion-sensitivity brackets and
      timestep sensitivity are reported as their own activity, distinct from
      comparison against observation.

\item \textbf{One narrow observational question with an answer either way.}
      Whether a westward-propagating barotropic branch is separable at all in a
      space--time spectral decomposition of reanalysis geopotential height
      \citep{hersbach2020,kalnay1996}, at which level, and with the spread between
      two independent reanalyses carried as observational uncertainty. If no such
      branch is resolvable, that is itself a finding about which level the
      comparison belongs at.
\end{enumerate}

\subsection{What is not claimed}

Stating this plainly is better than having a referee state it. No new equation,
mechanism or wave type is proposed. The numerical framework is Dedalus
\citep{burns2020} on the spherical basis of \citet{vasil2019}, used as documented.
This is not the first spherical barotropic-instability eigenvalue calculation
\citep{skiba2004,skiba2008}; not the first quantification of divergent
Rossby-mode dispersion \citep{longuethiggins1968,kasahara1976}; not a new
observational diagnostic, since both the space--time decomposition and the
correction of an observed ground-relative phase speed for advection by the mean
flow before comparison with an intrinsic prediction are standard; and not the
counter-propagating-wave picture, which is \citet{bretherton1966}'s, named by
\citet{hoskins1985} and presented by \citet{heifetz1999} explicitly as a
pedagogical framework.

Three limitations are stated here rather than deferred. The stability analysis is
posed for the \emph{nondivergent} barotropic potential-vorticity equation, while
the nonlinear runs integrate the divergent shallow-water system whose deformation
radius is comparable to the jet width; growth rates are therefore nondivergent
modal growth rates, and \citet{ripa1983} gives the general stability conditions a
fully divergent treatment would need. A normal-mode spectrum is not a stability
proof: the linearised operator is non-normal, so an all-real spectrum establishes
asymptotic stability to infinitesimal normal-mode disturbances and does not
exclude finite-time transient growth. And the model carries one degree of freedom
in the vertical, so it cannot represent baroclinic conversion --- which matters
directly for the observational comparison, since the observed extratropical
troposphere is dominated by eastward-propagating baroclinic disturbances that this
system does not describe.

\subsection{Organisation}

Section~\ref{sec:related} places the work in the literature by theory area.
Section~\ref{sec:theory} derives the framework from
$\mathrm{D}q/\mathrm{D}t = 0$, in both regimes. Section~\ref{sec:methods} describes
the solver, the resolution ladder and the verification protocol.
Section~\ref{sec:results} reports the wave campaign, the instability campaign and
the observational closure. Section~\ref{sec:discussion} returns to the planetary
regime question the rotation sweep opens.
